\ysrc{191}
这种相容性涉及模曲线的完备（上）同调，其细节留待下文 §3.2.11 讨论；眼下只须指出：模曲线 $Y(N)$ 的上同调 $H^1(Y(N),\Q_p)$ 中出现的每一组 Hecke 特征值，都对应 $\Gal_\Q$ 在某个 $\Q_p$-向量空间上的一个连续二维表示 $\rho$。$\rho$ 到素数 $\ell\neq p$ 处分解群上的限制 $\rho|_{\Gal_{\Q_\ell}}$，完全由 $\ell$ 处 Hecke 算子的特征值决定（至少差一个 Frobenius 半单性，而后者预期自动成立）。但 $p$ 处 Hecke 算子的作用只决定与 de Rham 表示 $\rho|_{\Gal_{\Q_p}}$ 相伴的 Weil--Deligne 表示。决定 $\rho|_{\Gal_{\Q_p}}$ 时所缺的那一块是 Hodge 滤过：它神秘地由远离 $p$ 处的 Hecke 作用决定，却不由 $p$ 处的 Hecke 作用捕捉。

于是 Breuil 最初的想法之一是：涉及完备上同调与 $p$-进局部 Langlands 对应的 $p$-进局部--整体相容性，应当使人能从完备上同调 $\widehat H^1(Y(N),\Q_p)$（定义见下文 (3.2.12)）中相应 Hecke 特征空间上的 $\mathrm{GL}_2(\Q_p)$-表示读出 Hodge 滤过，而不必用到远离 $p$ 处 Hecke 算子的作用。事实上，Hodge 滤过应当由该特征空间在 $\widehat H^1(Y(N),\Q_p)$ 中的局部解析向量决定。这一想法的推广，例如可见于 Breuil 的基座（socle）猜想（见 §9.6.33）：Hodge 滤过相对于 Frobenius 稳定旗的相对位置，由某个局部解析表示的基座决定。

把一个（经典）$(\varphi,\Gamma)$-模 $D$ 换成它在 Robba 环 $\mathcal R$ 上的变体 $D_{\mathrm{rig}}$，好处在于：即使 $D$ 不可约，该 $(\varphi,\Gamma)$-模也可能变得可约。例如当 $D$ 来自晶体表示或半稳定表示（更一般地，来自 crystabelline 或 semistabelline 表示）时，就会出现这种情况。完全可约的 $(\varphi,\Gamma)$-模称为 trianguline，Colmez [Col08] 曾对它们作过详尽研究。就 $p$-进局部 Langlands 对应而言，表示 $\rho$ 是 trianguline 的（即 $D_{\mathrm{rig}}(\rho)$ 是 trianguline 的）当且仅当局部解析向量 $\Pi(\rho)^{\mathrm{an}}$ 含有（局部解析主级数表示的）不可约成分。此时 $\Pi(\rho)^{\mathrm{an}}$ 通常是两个抛物诱导表示的扩张，而这些表示的扩张类由 $D_{\mathrm{rig}}(D)$ 所定义的秩 $1$ $(\varphi,\Gamma)$-模的扩张类决定（反之亦然）。特别地，当 $\rho$ 是晶体表示或半稳定表示时，这说明 Weil--Deligne 表示 $\mathrm{WD}(\rho)$ 上的 Hodge 滤过由局部解析向量 $\Pi(\rho)^{\mathrm{an}}$ 决定，与上述讨论的预期一致。

\subsubsection*{3.2.7.\quad $\mathrm{GL}_2(\Q_p)$ 的模 $p$ 局部 Langlands 对应}

到此为止的讨论中，我们把 $p$-进局部 Langlands 对应放在 $\mathrm{GL}_2(\Q_p)$ 的 $p$-进表示的框架里。然而，模 $p$ 的考量从这一理论肇始就起了根本作用。Breuil 在 [Bre03a] 中（并基于 Barthel 与 Livné [BL94] 的早期工作）分类了 $\mathrm{GL}_2(\Q_p)$ 在有限域 $k$ 上的绝对不可约容许光滑表示，并大致（grosso modo）证明了它们形如
\[
\Bigl(c\text{-}\mathrm{Ind}_{\mathrm{GL}_2(\Z_p)Z}^{\mathrm{GL}_2(\Q_p)}\sigma\Bigr)\otimes_{\mathcal H(\sigma),x}k,
\]
其中 $Z$ 表示 $\mathrm{GL}_2(\Q_p)$ 的中心，$\mathcal H(\sigma)=\End\bigl(c\text{-}\mathrm{Ind}_{\mathrm{GL}_2(\Z_p)Z}^{\mathrm{GL}_2(\Q_p)}\sigma\bigr)$ 是相应的 Hecke 代数（它带有显式同构 $\mathcal H(\sigma)\cong k[S^{\pm1},T]$），而 $x:\mathcal H(\sigma)\to k$ 是 $\mathcal H(\sigma)$ 的某个给定特征。[Bre03a] 的关键贡献在于证明：在所谓的超奇异情形，即 $x(T)=0$ 的情形，这类表示不可约（正是这一论断为 $\mathrm{GL}_2(\Q_p)$ 的情形所特有）。

\ysrc{192}在上述分类的基础上，Breuil 得以定义模 $p$ 半单局部 Langlands 对应：给每个半单二维表示 $\overline{\rho}:\Gal_{\Q_p}\to\mathrm{GL}_2(k)$ 配上 $\mathrm{GL}_2(\Q_p)$ 在 $k$ 上的一个半单容许光滑表示 $\kappa(\overline{\rho})$。Breuil 进一步猜想，这一对应与所猜想的 $p$-进局部 Langlands 对应相容，意思是：若 $\rho:\Gal_{\Q_p}\to\mathrm{GL}_2(L)$ 连续，则应有
\[
\overline{\Pi(\rho)}^{\mathrm{ss}}\cong\kappa(\overline{\rho}^{\mathrm{ss}}).
\]
（此处所猜同构两端出现的记号 $(-)^{\mathrm{ss}}$ 表示：先把一个不变格模单值化元 $\varpi$ 取约化，再半单化。又注意，两端的约化都只在半单化之后才唯一确定，因此把 $p$-进情形与模 $p$ 情形作此类比较时，过渡到半单化是自然的。）

随后 Berger [Ber10] 用 $(\varphi,\Gamma)$-模给出了 Breuil 模 $p$ 对应的一种描述，并借此在 [BB10] 的语境中验证了 Breuil 所猜的相容性。同样的 $(\varphi,\Gamma)$-模论证可用来证明：这一相容性自然地推广到 [Col10b] 的一般 $p$-进局部 Langlands 对应。

Colmez 随后定义了完整的模 $p$ 局部 Langlands 对应，即把该对应定义成如下样子：当 $\overline{\rho}$ 可约但不可分解时，它会把相应的特征扩张类也考虑进去。更确切地说，对每个表示 $\overline{\rho}:\Gal_{\Q_p}\to\mathrm{GL}_2(k)$，都可以定义一个相伴表示 $\kappa(\overline{\rho})$，使得 $V(\kappa(\overline{\rho}))=\overline{\rho}$。若 $\overline{\rho}$ 不是模 $p$ 分圆特征经平凡特征扩张所得的扭，则 $\kappa(\overline{\rho})$ 事实上可由这最后一个性质连同下述条件刻画：它没有有限维的 $\mathrm{GL}_2(\Q_p)$-不变子对象或商对象。（例如见 [Eme11a, Thm. 3.3.2]。）

在例外情形，即 $\overline{\rho}$ 是模 $p$ 分圆特征经平凡特征扩张所得的扭时，$\kappa(\overline{\rho})$ 的「正确」定义（我们指与 §7.2 所解释的对应相容的那个定义）与半单化不相容。事实上，即便在此情形，也可以改动 $\kappa(\overline{\rho})$ 的定义使之相容，Colmez 确实采用了这一替代定义；但这一替代的 $\kappa(\overline{\rho})$ 不满足与形变的相容性，其形式详见下文。（Colmez 的定义与我们的定义之间的差别，更精确的描述见 [CEG$^+$18, Rem. 6.22]。）

\subsubsection*{3.2.8.\quad 形变空间上的 $p$-进局部 Langlands 与 Breuil--Mézard 猜想的证明}

如前所述，$\mathrm{GL}_2(\Q_p)$ 的 $p$-进局部 Langlands 对应的构造，与「描述该对应如何在伽罗瓦表示的形变空间上变化」密切相关；下面转向讨论这一描述在文献中采取的各种形式。所有这些描述的第一个关键输入是：Colmez 的函子 $V$ 在模 $p$ 与 $p$-进整层面都有良好定义，因而可用于在形变论框架中研究 $p$-进局部 Langlands 对应。

[Eme11a] 对表示 $\overline{\rho}:\Gal_{\Q_p}\to\mathrm{GL}_2(k)$ 的形变给出了该对应的一种形变论描述，其中 $\overline{\rho}$ 不是分圆特征经平凡特征扩张所得的扭。下文为简明起见，假定 $\overline{\rho}$ 也只有平凡自同态，于是可以并且确实取 $R_p$ 为 $\overline{\rho}$ 带某个固定行列式的泛形变 $\mathcal O$-代数（而不必与加框架形变打交道）。为叙述简便，再回到先前的假定，即该固定行列式为 $\varepsilon^{-1}$；相应地，现在假定所考虑的一切 $\mathrm{GL}_2$ 表示都有平凡中心特征，从而把它们看作 $\mathrm{PGL}_2$ 的表示。

\ysrc{193}此时，Colmez [Col10b] 与 Kisin [Kis10a] 的论证可概括如下：存在 $\mathrm{PGL}_2(\Q_p)$ 的一个表示 $\pi^{\mathrm{univ}}$，作用在 $R_p$-模上，该模既是 $\mathfrak m_{R_p}$-进完备的、又是可正交归一化（orthonormalizable）的，它形变光滑表示 $\kappa(\overline{\rho})$，并具有性质 $V(\pi^{\mathrm{univ}})\cong\overline{\rho}^{\mathrm{univ}}$，即 $\overline{\rho}$ 的泛形变。于是，$p$-进局部 Langlands 对应配给 $\overline{\rho}$ 的任一具体形变的 $\mathrm{PGL}_2(\Q_p)$-表示，就通过沿该形变对 $\pi^{\mathrm{univ}}$ 作特殊化得到。

在 [Paš13] 中，给出了该对应的一种形变论描述，它比刚才回顾的表述更灵活、也更具概念性。为使描述尽可能简单，假定 $\overline{\rho}$ 不是平凡特征经分圆特征扩张所得的扭。（注意这与上述例外情形不同；被排除的那种扩张把两个特征以相反次序排列。）于是令 $\widetilde P$ 表示 Pontryagin 对偶 $\kappa(\overline{\rho})^{\vee}$ 在 $\mathrm{PGL}_2(\Q_p)$ 于紧 $\mathcal O$-模上的某个适当表示范畴中的投射包络（projective envelope）。用上标 $d$ 表示伪紧无 $\mathcal O$-挠 $\mathcal O$-模上的 Schikhof 对偶，即 $\Hom^{\mathrm{cont}}_{\mathcal O}(-,\mathcal O)$。[Paš13] 的主要结果之一是存在典范同构
\begin{equation}\tag{3.2.9}\label{eq:a4-3-2-9}
R_p\cong\End(\widetilde P),
\end{equation}
于是 $\widetilde P$ 可自然地看作「生活在 $\Spf R_p$ 之上」，而 $V(\widetilde P^{\,d})$ 上随之而来的 $R_p$-结构把 $V(\widetilde P^{\,d})$ 实现为 $\overline{\rho}$ 的泛形变 $r^{\mathrm{univ}}$。

若 $\rho:\Gal_{\Q_p}\to\mathrm{GL}_2(L)$ 有一个形变 $\overline{\rho}$ 的格，对应某个态射 $x:R_p\to\mathcal O$，则酉 Banach 空间表示 $\Pi(\rho)$ 由下述同构决定：
\[
\Pi(\rho)\cong L\otimes_{\mathcal O}\bigl(\widetilde P\otimes_{R_p,x}\mathcal O\bigr)^{d}.
\]
若 $\overline{\rho}$ 不是分圆特征经平凡特征扩张所得的扭（从而上述两个例外情形都不出现），则 [Paš13] 进一步证明 $\widetilde P$ 在 $R_p$ 上拓扑平坦（意思是函子 $-\widehat\otimes_{R_p}\widetilde P$ 在伪紧 $R_p$-模上正合）。此时，如上配给 $\overline{\rho}$ 的表示 $\pi^{\mathrm{univ}}$ 与 $\widetilde P$ 由同构
\[
\pi^{\mathrm{univ}}\cong\Hom^{\mathrm{cont}}_{R_p}(\widetilde P,R_p)
\]
相联系（这是 $R_p[\mathrm{PGL}_2(\Q_p)]$-模的同构）。

在 [Kis09a] 中，Kisin 给出了 $p$-进局部 Langlands 对应的一种颇为不同的描述：它不是在整个 $\Spf R_p$ 上，而是在参数化某个固定局部代数型 $\sigma$ 的提升的轨迹 $\Spf R_p(\sigma)$ 上。对于由 $\Spf R_p(\sigma)$ 参数化的 $\overline{\rho}$ 的形变 $\rho$，Colmez 对 Breuil 猜想的证明——该猜想把 $\Pi(\rho)$ 描述为局部代数表示 $\Pi(\rho)_{\mathrm{alg}}$ 的某个完备化——表明：$p$-进局部 Langlands 对应限制到 $\Spf R_p(\sigma)[1/p]$ 上所得者（记作 $\pi(\sigma)$）可以描述为 $c\text{-}\mathrm{Ind}_{\mathrm{PGL}_2(\Z_p)}^{\mathrm{PGL}_2(\Q_p)}\sigma$ 的某个完备化。

\ysrc{194}选定某个 $\mathrm{PGL}_2(\Z_p)$-不变格 $\sigma^{\circ}\subset\sigma$，它在约化 $\overline{\sigma}$ 中的 Jordan--Hölder 滤过决定了 $\overline{\pi}(\sigma)$ 上的相应滤过（即 $\pi(\sigma)$ 中某个不变格模 $\varpi$ 的约化），其各项为形如 $c\text{-}\mathrm{Ind}_{\mathrm{PGL}_2(\Z_p)}^{\mathrm{PGL}_2(\Q_p)}\overline{\sigma}_i$ 的表示的完备化，其中 $\overline{\sigma}_i$ 遍历 $\overline{\sigma}$ 的 Jordan--Hölder 因子。对每个 $\overline{\sigma}_i$，令 $\sigma_i$ 是 $\mathrm{PGL}_2(\Q_p)$ 的一个代数 $L$-表示，使 $\sigma_i^{\circ}\otimes_{\mathcal O}k\cong\overline{\sigma}_i$，并置 $R_p(\overline{\sigma}_i)=R_p(\sigma_i)\otimes_{\mathcal O}k$。现在，$\mathrm{PGL}_2(\Q_p)$ 不可约表示的显式分类，加上 Colmez 函子 $V$ 与约化相容这一事实，表明：能贡献非零完备化的 $\overline{\sigma}_i$ 只能是那些作为 $\overline{\rho}$ 的 Serre 权出现的；并且若某个 Serre 权 $\overline{\sigma}_i$ 贡献一个完备化，则该完备化必为 $c\text{-}\mathrm{Ind}_{\mathrm{PGL}_2(\Z_p)}^{\mathrm{PGL}_2(\Q_p)}\overline{\sigma}_i$ 的 $\mathfrak m_x$-进完备化，其中 $\mathfrak m_x$ 是特征 $x:\mathcal H(\overline{\sigma}_i)\to k$ 的核，而该特征使 $\bigl(c\text{-}\mathrm{Ind}_{\mathrm{PGL}_2(\Z_p)}^{\mathrm{PGL}_2(\Q_p)}\overline{\sigma}_i\bigr)\otimes_{\mathcal H(\sigma),x}k$ 成为 $\kappa(\overline{\rho})$ 的一个成分。此外，这个完备化若看作 $R_p(\sigma)$-模，实际上支撑在 $R_p(\overline{\sigma}_i)$ 上，而 $R_p(\overline{\sigma}_i)$ 的 $\mathfrak m_x$-进完备化被显式地等同于 $\mathcal H(\overline{\sigma}_i)$ 的 $\mathfrak m_x$-进完备化。这样就得到 Breuil--Mézard 型的不等式，但其方向与上面 3.2.3 所述论证给出的不等式相反。把这两个不等式结合起来便得到等式，即 Breuil--Mézard 猜想。如 3.2.3 所指出的，该猜想随后蕴含每个拼接模 $M_\infty(\sigma^{\circ})$ 在环 $R_p(\sigma)$ 上忠实，从而给出型为 $\sigma$ 的整体形变的模性提升定理。这就是 [Kis09a] 所采取的证明 Fontaine--Mazur 猜想的途径。

在上一段勾勒的论证中，之所以只得到不等式（而非 Breuil--Mézard 猜想的精确等式），多少是因为完备化（作为射影极限的一种形式）先验地并不右正合，因此先验地并不清楚：$\overline{\pi}(\sigma)$ 中可能出现的 $c\text{-}\mathrm{Ind}_{\mathrm{PGL}_2(\Z_p)}^{\mathrm{PGL}_2(\Q_p)}\overline{\sigma}_i$ 的一切可能完备化，是否真的都出现。

这类正合性问题归根结底属同调本性，因此，手头有更强的同调信息便能使论证加强，这或许并不令人意外。Paškūnas 在 [Paš15] 中做到了这一点。利用上述 $\widetilde P$ 的构造，连同 Colmez 关于局部代数向量的结果，可以发现 $R_p(\sigma)$ 恰好是 $R_p$-模 $\Hom_{\mathrm{PGL}_2(\Z_p)}(\sigma^{\circ},\widetilde P^{\,\vee})$ 的支撑（无论怎样选取 $\sigma$ 中的不变格 $\sigma^{\circ}$）。事实上，更方便的是与 $R_p$-模
\begin{equation}\tag{3.2.10}\label{eq:a4-3-2-10}
M(\sigma^{\circ}):=\widetilde P\otimes_{\mathcal O[[\mathrm{PGL}_2(\Z_p)]]}\sigma^{\circ}
\end{equation}
打交道，因为可以证明它是凝聚 $R_p$-模。由于 $\widetilde P$ 作为 $\mathrm{PGL}_2(\Q_p)$-表示是投射的，它特别地作为 $\mathrm{PGL}_2(\Z_p)$-表示也是投射的，故由上述公式隐含定义的函子 $M(-)$ 正合。用这个函子作论证，基本上可以把 Kisin 论证的两支合并成单一论证，直接得到 Breuil--Mézard 等式。

\subsubsection*{3.2.11.\quad 局部--整体相容性}

[Eme11a] 证明了一条局部--整体相容性，把 $p$-进局部 Langlands 对应与完备上同调联系起来；它是经典局部--整体相容性 [Lan73, Del73, Car83] 的 $p$-进类比，后者把无穷层（infinite level）的上同调与经典局部 Langlands 对应联系起来。

\ysrc{195}为了陈述该结果，像上面那样固定一个模不可约剩余表示 $r:\Gal_\Q\to\mathrm{GL}_2(k)$，并置 $\rho:=r|_{\Gal_{\Q_p}}$；在驯顺级 $N$ 上工作，其中 $N$ 被 $r$ 的与 $p$ 互素的 Artin 导子整除。与完备同调 (3.1.5) 类似，令
\[
\widehat H^1:=\widehat H^1_{\text{ét}}\bigl(Y(N),\mathcal O\bigr)
\]
表示驯顺级 $N$ 的「模曲线」（同样，更确切地说，是 $\mathrm{PGL}_2$-对称空间的同余商，它们是 $\mathrm{GL}_2$ 的通常模曲线各连通分支之并）的 $p$-进完备平展上同调；按定义
\begin{equation}\tag{3.2.12}\label{eq:a4-3-2-12}
\widehat H^1:=\varprojlim_{s}\ \varinjlim_{n}\ H^1_{\text{ét}}\bigl(Y(Np^n),\mathcal O/\varpi^s\bigr),
\end{equation}
而 $\widehat H^1$ 自然地带有群 $\Gal_\Q$ 与 $\mathrm{PGL}_2(\Q_p)$ 的交换作用，以及 Hecke 代数 $\mathbb T$ 的作用（$\mathbb T$ 由素数 $l\nmid Np$ 处的通常 Hecke 算子 $S_l,T_l$ 生成）。表示 $r$ 给出 $\mathbb T$ 的一个极大理想 $\mathfrak m$，于是可作相应的局部化 $\widehat H^1_{\mathfrak m}$。$\overline{r}$ 在 $\mathbb T_{\mathfrak m}$ 上也有泛 pro-模形变 $\rho^{\mathrm{univ}}$；若置 $\rho:=r|_{\Gal_{\Q_p}}$，则由 $\Gal_\Q$ 到 $\Gal_{\Q_p}$ 的限制诱导出态射 $R_p\to\mathbb T_{\mathfrak m}$。

现在假定 $\rho$ 不是分圆特征经平凡特征扩张所得的扭；[Eme11a] 的主要结果称存在同构（$\mathbb T_{\mathfrak m}[\Gal_\Q\times\mathrm{PGL}_2(\Q_p)]$-模的同构）
\begin{equation}\tag{3.2.13}\label{eq:a4-3-2-13}
\widehat H^1_{\mathfrak m}\cong(\rho^{\mathrm{univ}}\otimes_{R_p}\pi^{\mathrm{univ}})\,\widehat{\otimes}_{R_p}\,\mathcal O,
\end{equation}
其中 $\pi^{\mathrm{univ}}$ 是 §3.2.8 讨论过的 $R_p$ 上可正交归一化的 $\mathrm{PGL}_2(\Q_p)$-表示，$\mathcal O$ 是 $\mathbb T_{\mathfrak m}$ 的一个无 $p$-挠、余有限生成模\footnote{即某个有限生成 $\mathbb T_{\mathfrak m}$-模的 $\mathcal O$-对偶。}，由与所选驯顺级相关的「旧形式」（oldforms）构成（通过映射 $R_p\to\mathbb T_{\mathfrak m}$ 看作 $R_p$-模）。\footnote{若在满级 $N$ 上工作，再取遍所有级 $N$ 的余极限，则所得的旧形式极限空间带有 $\mathrm{GL}_2(\mathbb A^p_f)$ 的作用，并可用 [EH14] 与 [HM18] 的族化局部 Langlands 对应来描述。这与第 8 节简要讨论的范畴局部 Langlands 的 $\ell\neq p$ 情形有关，并在整体猜想（下文猜想 9.3.2）中起关键作用。不过为叙述简便，此处不再多说。}若 $\rho$ 进一步不是平凡特征经分圆特征扩张所得的扭，则过渡到连续 $\mathcal O$-对偶，并顾及此情形下 $\pi^{\mathrm{univ}}$ 与 $\widetilde P$ 的关系，可把 (3.2.13) 重述为同构
\begin{equation}\tag{3.2.14}\label{eq:a4-3-2-14}
(\widehat H_1)_{\mathfrak m}\cong\widetilde P\,\widehat{\otimes}_{R_p}(\rho^{\mathrm{univ}})^{\vee}\otimes_{R_p}\mathcal O^{*},
\end{equation}
其中 $\widehat H_1=\Hom_{\mathcal O}(\widehat H^1,\mathcal O)$ 表示驯顺级 $N$ 的模曲线的完备同调（如 (3.1.5) 那样），$(\rho^{\mathrm{univ}})^{\vee}$ 表示 $\rho^{\mathrm{univ}}$ 的 $R_p$-对偶，$\mathcal O^{*}$ 表示 $\mathcal O$ 的 $\mathcal O$-对偶（它现在是有限生成 $\mathbb T_{\mathfrak m}$-模）。该同构的证明结合了两样东西：一个是 Morita 理论式的论证，类似于 [Paš13] 中发展出的更精细的此类论证；另一个是插值论证，与 [CDP14] 中出现的「捕获」（capture）论证相关。关键输入是经典局部--整体相容性与 Berger--Breuil 的结果（当 $\rho$ 是 crystabelline 时，$\Pi(\rho)$ 是 $\Pi(\rho)_{\mathrm{alg}}$ 的唯一酉完备化）的结合：经典局部--整体相容性给出 $\widehat H^1$ 中若干份 $\Pi(\rho)_{\mathrm{alg}}$ 的拷贝，而 Berger--Breuil 的结果表明这些拷贝在 $\widehat H^1$ 中的闭包必同构于 $\Pi(\rho)$。

\ysrc{196}局部--整体相容性给出 Fontaine--Mazur 猜想的另一条途径。即：若 $x:\mathbb T_{\mathfrak m}\to L$ 是相应于 $\overline{r}$ 的某个 pro-模形变 $r$ 的同态，其核为 $\mathfrak p_x$，且置 $\rho:=r|_{\Gal_{\Q_p}}$，则在局部--整体相容性同构 (3.2.14) 中过渡到 $\mathfrak p_x$-挠，便得同构
\begin{equation}\tag{3.2.15}\label{eq:a4-3-2-15}
\widehat H^1[\mathfrak p_x]\otimes_{\mathcal O}E\cong r\otimes\Pi(\rho)
\end{equation}
（此处略去了旧形式的贡献；若取驯顺级与 $r$ 的与 $p$ 互素的导子一致，则它们的贡献确实是平凡的）。若 $\rho$ 是 Hodge--Tate 权互不相同的潜在半稳定表示，从而由 Colmez 的结果 $\Pi(\rho)_{\mathrm{alg}}\neq0$，则可断定由 $x$ 给出的 Hecke 特征值组出现在 $\widehat H^1_{\mathrm{alg}}$ 中。再把完备上同调中的局部代数向量与模曲线上局部系统的经典上同调作比较，便知 $x$ 从而 $r$ 来自一个经典模形式。

因此，关于 $\overline{r}$ 的形变的 Fontaine--Mazur 猜想归约为证明：所有这样的形变都是 pro-模的。在 Taylor--Wiles 方法所需的对 $\overline{r}$ 的标准假设下，这可以通过证明一个「大 $R=\mathbb T$」定理来完成，即把 $\mathbb T_{\mathfrak m}$ 等同于 $\overline{r}$ 的、在 $p$ 外极小分歧的形变的泛形变 $\mathcal O$-代数。这类定理最早在 [Böc01] 中证明，所用的是 Gouvêa 与 Mazur 的无限蕨（infinite fern），连同小权情形的模性提升定理（此时形变环 $R_p(\sigma)$ 可借助 Fontaine--Laffaille 理论显式描述）。此外我们指出：(3.2.15) 与 $p$-进 Langlands 对应的构造蕴含，对每个 $x:\mathbb T_{\mathfrak m}\to L$，$\widehat H^1[\mathfrak p_x]\otimes_{\mathcal O}L$ 非零；这比「相伴形变 $r$ 是 pro-模的」这一事实先验地更强。

因此 [Eme11a] 的结果给出了 [Kis09a] 诸结果的另一证明（假设略有不同）；两个证明都依赖 Taylor--Wiles--Kisin 拼接与 $p$-进 Langlands 对应，其区别在于：[Kis09a] 的进路是用 $p$-进 Langlands 对应对环 $R_p(\sigma)$ 证明足够多的东西，以推出拼接 $R_\infty(\sigma)$-模 $M_\infty(\sigma^{\circ})$ 忠实；而 [Eme11a] 的进路可看作是用 $p$-进 Langlands 对应来计算函子 $M_\infty$。

把 [Eme11a] 主要结果重述为关于完备同调的 (3.2.14)，这极大地暗示了两件事。第一，从本讲义的视角看（我们把伽罗瓦形变环换成 L-参数的模叠），它暗示可以用如下方式计算模曲线的完备上同调：把 $\mathrm{PGL}_2(\Q_p)$-表示的某个典范层，从局部参数的叠拉回到整体参数的叠上；我们在第 9 节追循这一提示。第二，它暗示可以用 Taylor--Wiles 拼接方法从完备同调 $\widehat H_1$ 中提取 $\widetilde P$；这一可能已在 [CEG$^+$18] 中实现。特别地，可以发现拼接函子 $M_\infty$「就是」由 (3.2.10) 隐含描述的函子 $M$；确切地说，它等于 $M$ 沿态射 $R_p\to R_\infty$ 的标量扩张。（不过注意：用这种方式，或用任何别的整体构造来构造 §7.2 的层 $\mathcal L_\infty$，看来毫无希望。）

\ysrc{197}
\subsubsection*{3.2.16.\quad 与经典局部 Langlands 证明的比较}

我们借整体手段构造 $p$-进局部 Langlands 对应，这可以与经典局部 Langlands 对应的构造相比较——后者同样通过把局部表示整体地实现来展开 [HT01, Hen00, Sch13]。然而，如 [CEG$^+$16, §1.3] 所论，我们的构造与经典构造之间有若干显著差别。经典上，人们先用 Weil--Deligne 表示的分类，以及 $\mathrm{GL}_n(\Q_p)$ 不可约光滑表示的平行分类 [BZ77, Zel80]，把问题归约为在不可约 $n$ 维 Weil--Deligne 表示与 $\mathrm{GL}_n(\Q_p)$ 的尖点表示之间构造双射。迹公式论证使人能把任一尖点表示（差一个扭）实现为某个尖点自守表示的局部成分，随后在某个志村簇的上同调中构造出相伴的 Weil--Deligne 表示（特别地，这一构造预设局部--整体相容性成立，正如 [CEG$^+$18] 的构造一样）。为验证如此构造的对应事实上是纯局部的，最终要依赖用 $L$-因子与 $\varepsilon$-因子对该对应所作的刻画。

在 $p$-进背景下，上述步骤大多不可用（至少目前如此）。例如，关于 $\mathrm{GL}_2(\Q_p)$ 容许酉 Banach 表示，已知的唯一分类是 [Paš13] 给出的，而它本身是用 $p$-进局部 Langlands 对应来表述的，并不独立于后者。此外，我们无法指望把局部 $p$-进伽罗瓦表示整体地实现，即便要求它是具有某些指定权与型的潜在半稳定表示也不行：这类表示的族是一维的（因而是不可数的），而 $\Gal_\Q$ 的模表示只有可数多个。这就是 Taylor--Wiles 拼接对 [CEG$^+$16, CEG$^+$18] 的进路如此关键的原因：在缺少 Bernstein--Zelevinsky 分类/迹公式论证（经典情形正是靠它把局部与整体语境联系起来）的 $p$-进情形，拼接提供了联系局部与整体语境的机制。

$p$-进与经典语境之间最大的差别或许是：在 $p$-进语境中，我们没有任何先验的刻画，能与经典情形用 $L$-因子与 $\varepsilon$-因子所作的刻画相类比；最接近的类比是要求 $p$-进对应与经典对应经「过渡到局部代数向量」而相容（在伽罗瓦一侧，这应平行于从潜在半稳定 Dieudonné 模过渡到其底层的 Weil--Deligne 表示）。然而，对群 $\mathrm{GL}_2(\Q_p)$，只有在主级数情形，这一要求才有可能提供足够信息以决定 $p$-进对应。这就是为什么在所有已知的 $\mathrm{GL}_2(\Q_p)$ 的 $p$-进局部 Langlands 对应证明中，（非分歧）主级数情形的考量如此关键。

\subsubsection*{3.2.17.\quad 拼接与解析框架}

在一系列论文 [BHS17b]、[BHS17a] 与 [BHS19] 中，作者们证明拼接构造不仅可用来证明关于自守形式与伽罗瓦表示的结果（例如证明存在某个自守形式，其相伴伽罗瓦表示具有指定的局部性质），还可用来证明关于超收敛 $p$-进自守形式及其相伴伽罗瓦表示的结果。

\ysrc{198}超收敛 $p$-进自守形式理论中的一个关键观察是：超收敛 $p$-进自守特征形式的 Hecke 特征值能良好地按族变化，并构成所谓的特征簇 $Y$。$p$-进自守形式本身（特征簇参数化的正是它们的 Hecke 特征值）在该特征簇上构成一个凝聚层 $\mathcal M$。特别地我们指出：在特征簇理论中，（超收敛有限斜率的）$p$-进自守形式空间可以描述为某个刚性解析空间上凝聚层的整体截面的（对偶），而该刚性解析空间与伽罗瓦表示形变空间的一般纤维密切相关。

用伽罗瓦表示纯然地刻画这个凝聚层 $\mathcal M$ 的支撑 $Y$，可视为 Fontaine--Mazur 猜想的超收敛版本：它相当于给出伽罗瓦表示 $\rho$ 的一个充分必要条件，保证 $\rho$ 与某个有限斜率的超收敛 $p$-进自守形式相伴。Kisin [Kis03] 构造了二维 $\Gal_\Q$-表示形变空间（与 $\GG_m$ 之积）的一个刚性解析子簇（即该文中的有限斜率空间），它用纯伽罗瓦论的语言定义（要求伽罗瓦表示在 $p$ 处 trianguline\footnote{trianguline 这个名字是 Colmez 后来才起的，他用的定义与 Kisin 文中的定义不同但等价。}、并在几乎处处非分歧），并猜想它等于 Coleman--Mazur 的特征曲线。这一构造事实上启发了 [BHS17b] 中 trianguline 簇的定义，以及下述 Fontaine--Mazur 猜想的超收敛版本（我们叙述得略不正式；$d=2$ 的情形最早由 Kisin 提出，见 [Kis03, 11.7(2), 11.8]）：一个不可约 $d$ 维奇 $p$-进整体伽罗瓦表示与 $\mathrm{GL}_d$ 上某个有限斜率的超收敛 $p$-进自守形式相伴，当且仅当它几乎处处非分歧、且在整除 $p$ 的位处 trianguline（回忆：$\rho:\Gal_\Q\to\mathrm{GL}_d(\Q_p)$ 是奇的，当且仅当 $\rho(c)$ 的迹为 $0$ 或 $\pm1$，其中 $c$ 表示复共轭）。（$d=2$ 的情形在辅助技术假设下于 [Eme11a, Thm. 1.2.4] 中得到证明；一个略弱的版本在略有不同的假设下于 [Kis09a, Cor., p. 642] 中建立。）

然而，这些构造与猜想只给出了凝聚层 $\mathcal M$ 的支撑 $Y$ 的刻画。对「用局部伽罗瓦表示以及伽罗瓦表示叠的整体到局部映射来纯然地刻画凝聚层 $\mathcal M$」这一精确猜想，我们参见 §9.6。认为这样的刻画应当存在的想法，源自如下洞见：拼接构造也能对特征簇实施。

更确切地说，正如 $R_\infty$ 上的 $M_\infty$（以及 $R_\infty(\sigma)$ 上的 $M_\infty(\sigma)$）经沿局部--整体映射的拉回与整体伽罗瓦形变环上的自守形式空间相联系，在（CM 域 $F$ 上满足拼接构造所需若干假设的定酉群 $G$ 的框架中）也存在特征簇 $X_p(\overline{\rho})$ 的拼接版本 [BHS17b]，以及 $X_p(\overline{\rho})$ 上的凝聚层 $\mathcal M_\infty$，使得 $\mathcal M$ 与 $Y$ 可经沿局部--整体映射的拉回从 $\mathcal M_\infty$ 与 $X_p(\overline{\rho})$ 重构出来。该构造的主要想法是模仿 M.E. [Eme06b] 构造特征簇的办法（用局部解析 Jacquet 函子），但用大的拼接模 $M_\infty$（相应地，其对偶 $\Pi_\infty$，即 $G(F_p)$ 的一个大 Banach 空间表示）代替 $p$-进自守形式空间。

\ysrc{199}[BHS17a] 与 [BHS19] 的关键想法是：用伽罗瓦论工具（即 §5.3 的空间 $\mathcal X_{G,\mathrm{tri}}$ 的几何）来分析空间 $X_p(\overline{\rho})$ 的局部几何，进而得到关于层 $\mathcal M_\infty$ 行为的信息。

\section{$(\varphi,\Gamma)$-模的模叠：Banach 情形}

现在我们回顾将要使用的 $(\varphi,\Gamma)$-模叠的（有时是猜想性的）构造。我们在两种不同框架中工作：本节处理 [EG23] 所论的「Banach」情形，它与 $G(\Q_p)$ 的光滑 $p$-进表示（猜想地）相关。第 5 节讨论「解析」情形，与 $G(\Q_p)$ 的局部解析表示相关。

本节非常简要地回顾 [EG23] 的一些主要结果。读者参阅 [EG24] 可获得对 [EG23] 的更详尽介绍与综述，参阅下文 §7.2 与 §7.6 可看到某些简单情形下这些叠的更显式描述。

下面考虑的叠是 [Eme] 意义下的形式代数叠（formal algebraic stack）；通俗地说，它们是「Artin 叠的形式类比物」，意思是它们容许形式概形（更确切地说，形式代数空间）的光滑覆盖。注意它们一般不是 $p$-进形式的；等价地说，它们的特殊（即模 $p$）纤维是真正的形式对象，而非代数对象。关于形式代数叠及其性质的基本综述，读者可参看 [EG23, App. A]。

\subsection{[EG23] 概述}

我们固定有限扩张 $F/\Q_p$，并使用作为 $\mathcal O$-代数的系数环，其中 $\mathcal O=\mathcal O_L$ 足够大。令 $\mathcal X_d$ 表示 $F$ 上秩 $d$ 的投射平展 $(\varphi,\Gamma)$-模的模叠（[EG23, Defn. 3.2.1] 意义下）。

\begin{rem}{4.1.1}
$\mathcal X_d$ 所参数化的对象有几种不同的可能描述。例如，可以把它们看作关于 $\Gal_F$ 的完全分圆扩张、或关于其 $\Z_p$ 子扩张的 $(\varphi,\Gamma)$-模；但有时把它们看作 $W(\C^{\flat})$ 上的平展 $(\Gal_F,\varphi)$-模（[EG23, Defn. 2.7.3] 意义下）也有帮助。它们还可以用（棱镜，prismatic）Laurent $F$-晶体来描述，见 [Min24]。
\end{rem}

按其定义本身（以及 Fontaine 的 $(\varphi,\Gamma)$-模理论 [Fon90]），$\mathcal X_d$ 的 $\overline{\F}_p$-点构成的群胚，与底层约化子叠 $\mathcal X_{d,\mathrm{red}}$ 的 $\overline{\F}_p$-点构成的群胚一致，自然地等价于连续表示 $\overline{\rho}:\Gal_F\to\mathrm{GL}_d(\overline{\F}_p)$ 的群胚。更一般地，若 $A$ 是任一有限 $\mathcal O$-代数，则群胚 $\mathcal X_d(A)$ 典范地等价于连续表示 $\Gal_F\to\mathrm{GL}_d(A)$ 的群胚。由此容易推出：泛伽罗瓦提升环是 $\mathcal X_d$ 在有限型点处（也就是在 $\overline{\F}_p$-点处；注意这些点自动可在 $\overline{\F}_p$ 的有限扩张上有代表）的通用（versal）环。因此可以把 $\mathcal X_d$ 看作 Mazur 伽罗瓦形变环的代数化。

\begin{rem}{4.1.2}
若 $A$ 不是有限 $\mathcal O$-代数（例如 $A=k[X]$），则通常没有与 $\mathcal X_d(A)$ 的某个对象相对应的伽罗瓦表示。存在 $\mathcal X_d$ 的一个极大子叠 $\mathcal X_d^{\mathrm{Gal}}$，在其上泛 $(\varphi,\Gamma)$-模来自 $\Gal_F$-表示；这个叠最初由 Wang--Erickson [WE18] 构造并研究。
\end{rem}

\ysrc{200}包含 $\mathcal X_d^{\mathrm{Gal}}\subset\mathcal X_d$ 在 $\overline{\F}_p$-点上诱导双射（且在这些点处是通用的），大体可以设想 $\mathcal X_d^{\mathrm{Gal}}$ 是 $\mathcal X_d$ 在其闭 $\overline{\F}_p$-点（这些点对应半单 $\Gal_F$-表示）处形式完备化之并。因此，尽管叠 $\mathcal X_d^{\mathrm{Gal}}$ 是「Langlands 参数」模叠的更显然的定义，叠 $\mathcal X_d$ 却有更丰富的几何结构；并且，例如从整体伽罗瓦表示之叠到 $\mathcal X_d$ 的限制映射虽经 $\mathcal X_d^{\mathrm{Gal}}$ 分解，我们仍预期：$\mathcal X_d$（更确切地说，它上面的拟凝聚层导出范畴）才是 $p$-进局部 Langlands 对应的自然舞台。

在 $p$-进 Langlands 纲领中，常常重要的是对所考虑的 $\Gal_F$-表示施加条件，例如要求它们是具有固定 Hodge--Tate 权的晶体表示，或更一般地，是某个固定惯性型的潜在晶体或半稳定表示。对任一 Hodge 型 $\lambda$ 与惯性型 $\tau$，都有相应的子叠 $\mathcal X_d^{\mathrm{crys},\lambda,\tau}$（以及类似的 $\mathcal X_d^{\mathrm{ss},\lambda,\tau}$），它由下述性质刻画：它是 $(\mathcal X_d)_{\mathcal O}$ 的唯一闭子叠，在 $\mathcal O$ 上平坦，且对任一有限平坦 $\mathcal O$-代数 $A$，其 $A$-值点群胚等价于满足相应 $p$-进 Hodge 论性质（即具有 Hodge 型 $\lambda$ 与惯性型 $\tau$ 的潜在晶体或半稳定）的 $A$-值 $\Gal_F$-表示群胚。下文中若 $\tau$ 平凡，则把它从记号中略去。\footnote{对下述构造，须假定 $L$「足够大」，例如使 $\tau$ 在 $L$ 上有定义。下文中我们大体上不提这一点。}

这些子叠在 [EG23] 中借助（除其他技巧外）Kisin 关于 Breuil--Kisin 模与晶体表示的结果 [Kis06] 构造；按 [BS23]，它们或许更自然地用（对数）棱镜 $F$-晶体来定义。

下述定理总结了 [EG23] 的主要结果。

\begin{thm}{4.1.3}\label{thm:a4-4-1-3}
\begin{enumerate}[label=(\arabic*),leftmargin=2.4em,itemsep=0.15em]
\item $\mathcal X_d$ 是 Noether 形式代数叠。
\item 每个子叠 $\mathcal X_d^{\mathrm{crys},\lambda,\tau}$ 与 $\mathcal X_d^{\mathrm{ss},\lambda,\tau}$ 都是 $p$-进形式代数叠。
\item $\mathcal X_d$ 的底层约化子叠 $\mathcal X_{d,\mathrm{red}}$（它是代数叠）在 $\overline{\F}_p$ 上有限型，且等维（equidimensional），维数为 $[F:\Q_p]d(d-1)/2$。
\item $\mathcal X_d^{\mathrm{crys},\lambda,\tau}$ 与 $\mathcal X_d^{\mathrm{ss},\lambda,\tau}$ 各自的特殊纤维（它是代数叠）等维，维数等于由 $\lambda$ 决定的旗簇的维数。特别地，若 $\lambda$ 正则，则这些叠各自等维，维数为 $[F:\Q_p]d(d-1)/2$。
\item $\mathcal X_{d,\mathrm{red}}$ 的不可约分支可按 Serre 权自然地标号。
\end{enumerate}
\end{thm}

\begin{rem}{4.1.4}
Serre 权的概念在 §6.1.38 中回顾。用 Serre 权给 $\mathcal X_{d,\mathrm{red}}$ 的不可约分支标号，其确切含义有些复杂；见 [EG23, §5.5]。粗略地说，描述如下。$\mathcal X_{d,\mathrm{red}}$ 的每个不可约分支都含有一个稠密开子叠，其 $\overline{\F}_p$-点对应 niveau 为 $1$ 的极大非分裂 $\Gal_F$-表示；也就是说，这些表示容许唯一按特征给出的滤过。这些特征的惯性权组成的有序组决定该不可约分支，而这些有序组对应 Serre 权，因为按定义 Serre 权是 $\mathrm{GL}_d(k)$ 的不可约 $\overline{\F}_p$-表示（的同构类），因而按其最高权编号。（作此辨认时有一个微妙之处：必须区分最高权模 $p-1$ 同余的 Serre 权；在伽罗瓦一侧，这表现为 peu ramifiée 表示与 très ramifiée 表示之间的区别。）
\end{rem}

\ysrc{201}
\begin{rem}{4.1.5}
定理 4.1.3 中 $[F:\Q_p]d(d-1)/2$ 出现两次——一次作为底层约化子叠 $\mathcal X_{d,\mathrm{red}}$ 的维数，一次作为正则 Hodge 型的潜在晶体/半稳定子叠特殊纤维的维数——并非巧合。它既是定理 4.1.3 证明中的重要成分，也支撑着 Breuil--Mézard 猜想（见 §6.1.38）。
\end{rem}

\begin{rem}{4.1.6}
据我们所知，文献中没有关于形式代数叠的一般维数理论。不过，认为这样的理论存在、且该维数可由通用环的维数算出，是合理的。若果真如此，则 [BIP23a] 的结果（它们特别计算了无限制局部伽罗瓦提升环的维数）蕴含 $\mathcal X_d$ 等维，维数为 $1+[F:\Q_p]d^2$。既然其底层约化子叠 $\mathcal X_{d,\mathrm{red}}$ 等维、维数为 $[F:\Q_p]d(d-1)/2$，可见 $\mathcal X_d$ 应有 $1+[F:\Q_p]d(d+1)/2$ 个纯形式的（包括 $p$-进方向在内的）方向。特别地，它（可证地）从不是 $p$-进形式代数叠。
\end{rem}

\begin{rem}{4.1.7}
同样据我们所知，文献中没有「形式代数叠是 lci（局部完全交）」这一说法的概念。我们预期合理的定义是存在的，且 $\mathcal X_d$ 是 lci；这同样应由 [BIP23a] 的结果推出（后者证明无限制局部伽罗瓦提升环是 lci）。（这一预期的一个版本由 Min 在 [Min24] 中实现，该文证明 $\mathcal X_d$ 的一个自然的导出变体实际上是经典的。）
\end{rem}

\subsection{固定行列式变体}

考虑这些叠的一个固定行列式的变体，往往很方便。若 $\chi:\Gal_F\to\mathcal O^{\times}$ 是特征，则令 $\mathcal X_d^{\chi}(A)$ 表示对 $(D,\theta)$ 的群胚，其中 $D$ 是系数在 $A$ 中、秩 $d$ 的投射平展 $(\varphi,\Gamma)$-模，$\theta$ 是 $\wedge^dD$ 与 $\chi$ 的一个等同。这些叠（及其潜在晶体与半稳定变体）在 [DEG] 的一篇附录中研究，定理 4.1.3 的类比在那里得到建立（唯一的差别是：在「Steinberg」Serre 权的情形，可以有多个相应的不可约分支，按 $\overline{\F}_p$ 中的 $d$ 次单位根编号）。特别地，$\mathcal X_d^{\chi}$ 的底层约化子叠同样等维，维数为 $[F:\Q_p]d(d-1)/2$。

\section{$(\varphi,\Gamma)$-模的模叠：解析情形}

与第 4 节类似，我们想在刚性解析空间的范畴上定义并研究 $(\varphi,\Gamma)$-模的叠。把这类对象看作 Fargues--Fontaine 曲线上的等变向量丛，有时更方便。本节大部分结果将出现在即将发表的论文 [HHS] 中。我们先回顾相应的框架，以及与 $(\varphi,\Gamma)$-模的比较。

\ysrc{202}
\subsection{Robba 环上的 $(\varphi,\Gamma)$-模}

回忆我们固定了 $\Q_p$ 的有限扩张 $F$ 及 $F$ 的一个代数闭包 $\overline{F}$。此外固定一个 $p^n$ 次单位根的相容系统 $\epsilon_{p^n}$，并记 $F_n=F(\epsilon_{p^n})$、$F_\infty=\bigcup_n F_n\subset\overline{F}$。

此外记 $\Gamma=\Gal(F_\infty/F)\supset\Gamma_n=\Gal(F_\infty/F_n)$。又注意 $\overline{F}$ 与 $F_\infty$ 的 $p$-进完备化都是 perfectoid 域。

联系于域 $F$（更确切地说，联系于分圆扩张 $F_\infty$ 与代数闭包 $\overline{F}$），可以定义 Fargues--Fontaine 曲线的两个版本。令 $\varpi^{\flat}\in\mathcal O_{\widehat F^{\flat}_\infty}$ 表示所选的一个伪单值化元（pseudo-uniformizer）。置
\begin{align*}
Y_{F_\infty}&=\Spa\bigl(W(\mathcal O_{\widehat F^{\flat}_\infty}),W(\mathcal O_{\widehat F^{\flat}_\infty})\bigr)\backslash V(p[\varpi^{\flat}]),\\
Y_{\overline{F}}&=\Spa\bigl(W(\mathcal O_{\widehat{\overline{F}}^{\flat}}),W(\mathcal O_{\widehat{\overline{F}}^{\flat}})\bigr)\backslash V(p[\varpi^{\flat}]).
\end{align*}
这些 $\Q_p$ 上的 adic 空间分别带有自同构 $\varphi$、以及 $\Gamma$（分别地 $\Gal_F$）的连续作用，且与 $\varphi$ 交换。Fargues--Fontaine 曲线 $X_{F_\infty}$ 与 $X_{\overline{F}}$ 是商
\[
X_{F_\infty}=Y_{F_\infty}/\varphi^{\Z},\qquad
X_{\overline{F}}=Y_{\overline{F}}/\varphi^{\Z},
\]
它们分别继承 $\Gamma$（分别地 $\Gal_F$）的连续作用。

注意曲线 $X_{F_\infty}$ 与 $X_{\overline{F}}$ 有一个由 Fontaine 通常的 $\theta$ 映射之核定义的特殊点 $\infty$，例如见 [FF18, 2.1.2]：
\begin{align*}
\theta_{F_\infty}&:W(\mathcal O_{\widehat F^{\flat}_\infty})\longrightarrow\mathcal O_{\widehat F_\infty},\\
\theta_{\overline{F}}&:W(\mathcal O_{\widehat{\overline{F}}^{\flat}})\longrightarrow\mathcal O_{\widehat{\overline{F}}}.
\end{align*}
由构造直接可见，点 $\infty$ 是 $\Gamma$（分别地 $\Gal_F$）作用的不动点。

我们定义下述（半）线性代数对象的范畴：
\begin{itemize}[leftmargin=2.2em,itemsep=0.15em]
\item $\Gamma$-等变向量丛在 $X_{F_\infty}$ 上的范畴 $\mathrm{VB}^{\Gamma}_{X_{F_\infty}}$；
\item $\Gal_F$-等变向量丛在 $X_{\overline{F}}$ 上的范畴 $\mathrm{VB}^{\Gal_F}_{X_{\overline{F}}}$；
\item $(\varphi,\Gamma)$-等变向量丛在 $Y_{F_\infty}$ 上的范畴 $\mathrm{VB}^{\varphi,\Gamma}_{Y_{F_\infty}}$；
\item $(\varphi,\Gal_F)$-等变向量丛在 $Y_{\overline{F}}$ 上的范畴 $\mathrm{VB}^{\varphi,\Gal_F}_{Y_{\overline{F}}}$。
\end{itemize}
事实上，上述所有范畴中都要求 $\Gamma$（分别地 $\Gal_F$）在向量丛上的作用按显然的意义连续（我们不明确写出）。由 Berger、Fargues--Fontaine 与 Kedlaya 的工作，已知这些各不相同的范畴都典范等价。

\begin{thm}{5.1.1}\label{thm:a4-5-1-1}
(i) 投影 $X_{\overline{F}}\to X_{F_\infty}$ 关于典范投影 $\Gal_F\to\Gamma$ 等变，并经拉回诱导范畴等价
\[
\mathrm{VB}^{\Gamma}_{X_{F_\infty}}\to\mathrm{VB}^{\Gal_F}_{X_{\overline{F}}}.
\]
对 $Y_{\overline{F}}\to Y_{F_\infty}$ 的相应陈述同样成立。\\
(ii) 投影 $Y_{F_\infty}\to X_{F_\infty}$ 是 $\Gamma$-等变的，并经拉回诱导范畴等价
\[
\mathrm{VB}^{\Gamma}_{X_{F_\infty}}\to\mathrm{VB}^{\varphi,\Gamma}_{Y_{F_\infty}}.
\]
对 $Y_{\overline{F}}\to X_{\overline{F}}$ 的相应陈述同样成立。
\end{thm}

% 本片止于印刷页 202（PDF 215）之末，定理 5.1.1 至此结束；续片自 PDF 216 页（印刷页 203）起。
